\chapter{Appendix A: Glosarium Istilah AI}
\label{chap:appendix_a}

\begin{description}
    \item[AI (Artificial Intelligence)] Kemampuan mesin untuk meniru fungsi kognitif manusia.
    \item[AGI (Artificial General Intelligence)] AI yang memiliki kecerdasan setara manusia di segala bidang.
    \item[Alignment] Proses memastikan AI bertindak sesuai dengan nilai-nilai manusia.
    \item[Bias] Prasangka atau kesalahan sistematis dalam data yang menyebabkan keputusan AI tidak adil.
    \item[Chatbot] Program komputer yang mensimulasikan percakapan manusia.
    \item[Chunking] Memecah teks panjang menjadi bagian-bagian kecil agar muat dalam Context Window.
    \item[Context Window] Batas memori jangka pendek LLM, diukur dalam token.
    \item[Deep Learning] Subset ML yang menggunakan jaringan saraf tiruan berlapis-lapis.
    \item[Embedding] Representasi numerik (vector) dari kata atau kalimat.
    \item[Fine-Tuning] Melatih ulang model AI yang sudah jadi dengan data spesifik untuk tugas tertentu.
    \item[Generative AI] AI yang bisa menciptakan konten baru (teks, gambar, audio).
    \item[Hallucination] Fenomena ketika AI memberikan informasi yang salah dengan sangat percaya diri.
    \item[Inference] Proses menjalankan model AI yang sudah dilatih untuk membuat prediksi.
    \item[LLM (Large Language Model)] Model AI yang dilatih dengan data teks dalam jumlah sangat besar.
    \item[Machine Learning] Kemampuan komputer belajar dari data tanpa diprogram eksplisit.
    \item[Multimodal] Kemampuan AI memproses berbagai jenis input (teks, gambar, suara) sekaligus.
    \item[NLP (Natural Language Processing)] Cabang AI yang berurusan dengan bahasa manusia.
    \item[Parameter] Variabel internal model yang dipelajari selama pelatihan (GPT-4 punya triliunan).
    \item[Pre-Training] Tahap awal pelatihan model dengan data mentah yang sangat besar.
    \item[Prompt] Input teks yang diberikan user kepada AI.
    \item[Prompt Engineering] Seni merancang prompt untuk mendapatkan output terbaik.
    \item[RAG (Retrieval Augmented Generation)] Teknik menggabungkan LLM dengan database eksternal untuk akurasi.
    \item[RLHF (Reinforcement Learning from Human Feedback)] Metode melatih AI dengan feedback manusia agar lebih aman dan bermanfaat.
    \item[Temperature] Parameter yang mengatur kreativitas/keacakan output AI (0.0 = deterministik, 1.0 = kreatif).
    \item[Token] Unit dasar teks yang diproses oleh LLM (bisa berupa kata atau bagian kata).
    \item[Transformer] Arsitektur neural network di balik semua LLM modern.
    \item[Vector Database] Database khusus untuk menyimpan embedding vector.
    \item[Zero-Shot] Meminta AI melakukan tugas tanpa memberikan contoh sebelumnya.
\end{description}

\chapter{Appendix B: Workbook Proyek AI (75+ Templat)}
\label{chap:appendix_b}

Bagian ini berisi kumpulan lembar kerja (worksheet) kosong yang dirancang untuk dicetak dan diisi. Gunakan untuk berbagai proyek AI Anda.

\foreach \n in {1,...,180}{
    \newpage
    \section*{Lembar Kerja Proyek AI - Halaman \#\n}

    \begin{tcolorbox}[colback=white, colframe=black, height=7cm, title=1. Identifikasi Masalah]
        \textbf{Masalah Utama:} \vspace{1.5cm} \hrule \vspace{0.5cm}
        \textbf{Target Pengguna:} \vspace{1.5cm}
    \end{tcolorbox}

    \vspace{0.5cm}

    \begin{tcolorbox}[colback=white, colframe=black, height=7cm, title=2. Solusi \& Data]
        \textbf{Solusi AI:} \vspace{1.5cm} \hrule \vspace{0.5cm}
        \textbf{Sumber Data:} \vspace{1.5cm}
    \end{tcolorbox}

    \vspace{0.5cm}

    \begin{tcolorbox}[colback=white, colframe=black, height=7cm, title=3. Catatan Eksperimen]
        \textbf{Prompt Awal:} \vspace{1.5cm} \hrule \vspace{0.5cm}
        \textbf{Hasil:} \vspace{1.5cm}
    \end{tcolorbox}
}

\chapter{Appendix C: Log Eksperimen Prompt}
\label{chap:appendix_c}

Gunakan tabel ini untuk mencatat iterasi prompt Anda. Ingat, prompt yang baik lahir dari revisi berulang-ulang.

\newpage
\section*{Log Iterasi Prompt}

\begin{center}
\begin{tabular}{|p{3cm}|p{8cm}|p{3cm}|}
\hline
\textbf{Versi} & \textbf{Prompt} & \textbf{Skor (1-10)} \\
\hline
v1.0 & & \\[2cm]
\hline
v1.1 & & \\[2cm]
\hline
v1.2 & & \\[2cm]
\hline
v1.3 & & \\[2cm]
\hline
v2.0 & & \\[2cm]
\hline
\end{tabular}
\end{center}

\vspace{1cm}

\begin{center}
\begin{tabular}{|p{3cm}|p{8cm}|p{3cm}|}
\hline
\textbf{Versi} & \textbf{Prompt} & \textbf{Skor (1-10)} \\
\hline
v1.0 & & \\[2cm]
\hline
v1.1 & & \\[2cm]
\hline
v1.2 & & \\[2cm]
\hline
v1.3 & & \\[2cm]
\hline
v2.0 & & \\[2cm]
\hline
\end{tabular}
\end{center}

\chapter{Appendix D: Cheat Sheet Pola Prompt}
\label{chap:appendix_d}

\begin{itemize}
    \item \textbf{Persona Pattern:} "Bertindaklah sebagai [Peran]..."
    \item \textbf{Audience Persona Pattern:} "Jelaskan ini kepada [Audiens]..."
    \item \textbf{Flipped Interaction Pattern:} "Tanyakan saya pertanyaan tentang [Topik] sampai Anda memiliki cukup informasi untuk melakukan [Tugas]."
    \item \textbf{Recipe Pattern:} "Saya ingin mencapai [Tujuan]. Saya tahu saya perlu melakukan [Langkah 1], [Langkah 2]. Tolong lengkapi langkah-langkahnya."
    \item \textbf{Alternative Approaches Pattern:} "Berikan 3 alternatif solusi untuk masalah ini dan bandingkan pro-kontranya."
    \item \textbf{Refinement Pattern:} "Kritik jawaban Anda sendiri dan berikan versi yang lebih baik."
\end{itemize}

\endinput
